% ============================================================
% Chapter 6: Experimental Design
%   Section: Evaluation Harness
% Included from chapters/06-experimental-design/experimental-design.tex
% Source: docs/decisions/evaluation.md; backend/src/evaluation/.
% ============================================================

\section{Evaluation Harness}
\label{sec:harness}

Every verdict-level number in this thesis is produced by one tool, the
evaluation harness (\path{backend/src/evaluation/}). It runs the
pipeline's own claim-checking stage --- the one behind the single-claim
endpoint, not a copy of it --- over a labelled set in X-Fact's format,
which both the X-Fact files and the custom set share, and scores the
results against the labels:

\begin{lstlisting}
uv run python -m src.evaluation.cli run --dataset <set> --model <model>
uv run python -m src.evaluation.cli report --run <run directory>
\end{lstlisting}

Two gaps in the pipeline were closed before it produced a single
number. The claim-checking stage took no language, so Spanish claims
would have been searched with the English word lists; and it took no
article context, so a custom-set claim could have been confirmed by
its own article. Both are now optional arguments: the harness passes
each row's language, and for the custom set the article's URL, which
the retriever excludes. X-Fact rows get no context: the harness gives
the pipeline nothing from the gold row that production would not have.

The model is an argument of the run and the provider is configuration
only, so the same command measures the local model and a hosted one.
The internal corpus is a snapshot copied into the run at its first
start, so a run does not change as the live store grows.

% ============================================================
% Chapter 6: Experimental Design
%   Section: Evaluation Harness
%     Subsection: Record Format and Run Layout
% Included from chapters/06-experimental-design/04-evaluation-harness/evaluation-harness.tex
% Source: docs/decisions/evaluation.md, "Record format", "Run layout,
% cache and resume"; record.py, store.py.
% ============================================================

\subsection{Record Format and Run Layout}
\label{subsec:harness-records}

\paragraph{One record per claim.} Each claim's result is one line of a
JSON Lines file, with every field present even on an error, so a report
never has to ask whether a key exists. A record holds the gold row's
fields, so a report needs nothing else; the verdict and confidence
after recalibration and the model's raw answer before it; the stage the
claim reached and whether its search or model failed; the queries sent,
with their kind; \emph{every} candidate source seen, once, with the
stage that cut it and why; the ranked evidence in the order the model
saw it, with each factor of its score, its stance, whether it was cited
and the quote; the full trace of progress events with timestamps; the
latency of each stage; the model calls with their tokens; and the
provenance --- model, provider, thresholds, harness version and code
commit. No scraped page body is stored.

\paragraph{Runs are keyed, so the same command resumes.} A run lives in
a directory named by a key derived from the dataset's hash, the model,
the effective thresholds, the scoring weights, the corpus mode and the
harness version: changing any of them starts a new run, while the same
command continues the old one. The commit is recorded on every record
but not keyed, which would discard a run on every commit. Editing a
label changes the dataset's hash, so two gold sets are never mixed in
one run.

\paragraph{Resuming.} Claims already scored are not run again; errors
are retried; a torn last line from a crash is cut and its claim re-run;
and claims whose search or model was unreachable can be re-run on
request, since those are outages, not results. Every record is flushed
and synced to disk as it is written. During the X-Fact pilot this rule
absorbed an interruption of the container runtime and two native
crashes, re-running exactly the claims each had lost
(Section~\ref{sec:results-pilot}). Each session of a run records its
own settings --- timeouts and concurrency --- a fix the pilot prompted.

\paragraph{What is committed.} Raw runs are not committed: they hold
third-party titles and quotes. The report of a run --- its metrics in
JSON, a readable report and a copy of the run's manifest --- is
committed, and those are the numbers this document quotes.


% ============================================================
% Chapter 6: Experimental Design
%   Section: Evaluation Harness
%     Subsection: Metrics and Confidence Intervals
% Included from chapters/06-experimental-design/04-evaluation-harness/evaluation-harness.tex
% Source: docs/decisions/evaluation.md, "What is reported apart",
% "Metrics"; metrics.py, report.py.
% ============================================================

\subsection{Metrics and Confidence Intervals}
\label{subsec:harness-metrics}

\paragraph{What is reported apart.} Each claim gets one outcome, the
first that applies: \emph{error} (the harness or the pipeline crashed),
\emph{search unavailable} (no web query was answered),
\emph{model unreachable}, or \emph{scored}. Only scored claims enter
the metrics. The others are listed by identifier and never counted as
errors of the system: on 25~September, 68 of 69 searches came back
empty, and folded into the error rate that would have measured the
search engine's uptime. On top of the outcome, two flags:

\begin{itemize}
    \item \textbf{Temporal leak}: a cited source was published after
    the claim (the annotation guide's Rule~3). Undetermined when the
    claim or the source has no date.
    \item \textbf{Verdict leak} (X-Fact only): a ranked source is on the
    fact-checker's own site, so the verdict may have been read rather
    than reached.
\end{itemize}

Every metric is reported on all scored claims and again without the
flagged ones. An answer the model gave in no usable format is scored as
the \texttt{UNVERIFIED} it became and counted per model.

\paragraph{Verdict metrics.} Over the five verdicts:

\begin{itemize}
    \item \textbf{accuracy}, and \textbf{macro-averaged $F_1$} over the
    classes present in the gold labels --- a class with no gold support
    is shown with its predictions but not averaged in --- with
    per-class precision, recall, $F_1$, support and predicted count, and
    the confusion matrix;
    \item \textbf{Cohen's $\kappa$} against the gold labels
    \cite{cohen1960coefficient}, the same statistic as the annotator's
    self-agreement, so that the system and the annotator are compared
    on one scale;
    \item \textbf{coverage} --- the share of claims given a definitive
    verdict rather than \texttt{UNVERIFIED} --- and \textbf{selective
    accuracy}, the accuracy on those claims, the standard pair for a
    classifier allowed to abstain \cite{geifman2017selective}.
\end{itemize}

\paragraph{Intervals.} Every statistic comes with a percentile 95\%
bootstrap interval over claims \cite{efron1993bootstrap}: 10,000
resamples, seeded (2026 by default) and with the seed in the report,
records ordered by identifier first so that the same run gives the same
interval however its claims finished. Two runs are compared with a
\emph{paired} bootstrap of their difference over the claims both
scored, alongside how many claims only one of them got right; a run
compared with itself gives exactly zero, which a test holds.

\paragraph{Breakdowns.} By language for both sets; by topic group,
claim type, source tier and origin (natural or synthetic) for the
custom set; by fact-checking site for X-Fact. Each with its size,
accuracy and interval, macro-$F_1$, coverage and selective accuracy.

\paragraph{Cost.} Each record carries its model calls, tokens and the
provider's answer time, measured inside the production client, so the
report gives the median and 90th-percentile latency per claim, tokens
per claim and per second, and the price per hundred claims from the
price table --- ``n/a'' without a price, never zero.


% ============================================================
% Chapter 6: Experimental Design
%   Section: Evaluation Harness
%     Subsection: Retrieval Evaluation
% Included from chapters/06-experimental-design/04-evaluation-harness/evaluation-harness.tex
% Source: docs/decisions/evaluation.md, "Retrieval evaluation";
% retrieval.py.
% ============================================================

\subsection{Retrieval Evaluation}
\label{subsec:harness-retrieval}

RQ2 asks what retrieval contributes. The harness answers in two ways.

\paragraph{Against the annotator's evidence.} For each claim with
reference links, it checks at three depths --- every candidate seen,
the ranked evidence the model was shown, and the evidence it cited ---
whether a reference \emph{link} was found, and whether a reference
\emph{domain} was; and it computes the reciprocal rank of the first hit
in the ranked list. Links are normalised with the retriever's own
function, imported rather than copied, so ``found'' means what the
pipeline means by it. For X-Fact, links on the fact-checker's own site
are removed from the references and counted, since they are the
verdict. A miss is not proof of failure --- another source can be as
good as the annotator's --- which is why domain hits are reported beside
link hits. The cut reasons of references that were found and then cut
show where they were lost.

\paragraph{Without references.} For any set: the share of claims with
any ranked evidence, candidates and ranked sources per claim, the share
cut at ranking, distinct and rated domains, the share from the internal
corpus, retrieval and ranking time, queries per claim, and which query
kind --- anchor, proposition or refutation --- found each ranked
source. All are means over claims with bootstrap intervals, by set and
by language.

\paragraph{Strategies.} A retrieval strategy is a run with its own
thresholds, corpus, engines or code, named by a label that enters the
run key. A retrieval-only mode stops before the model, so a strategy
costs only its search and fetch time, and a comparison command sets
strategies side by side, each paired against the first. This is what
the planned ablation (Table~\ref{tab:priorities}) uses: switching a
component off on the same frozen claims, in one window, rather than
comparing versions of the code that cannot be re-run on today's web.

\paragraph{Switching queries off.} Which of the three queries are sent
is a per-run threshold: a set of query kinds, all three by default.
Only when it is not all three does it enter the run key, so every run
made before the switch keeps its key.

\paragraph{Pooled relevance judgements.} The annotator's links are a
strict witness: another page can serve as well. So, as in TREC
\cite{voorhees2005trec}, every variant's top five ranked sources per
claim are pooled, a page found by two variants entering once; a person
judges each page once, in the terminal, as bearing on the claim or not;
and every variant is scored on the same judgements --- precision at
five, recall of the pool's relevant pages, reciprocal rank, and the
share of its top five that was judged. A variant pooled after the
judging adds only its own pages and reports how much of it was not
judged.


% ============================================================
% Chapter 6: Experimental Design
%   Section: Evaluation Harness
%     Subsection: Stage Attribution
% Included from chapters/06-experimental-design/04-evaluation-harness/evaluation-harness.tex
% Source: docs/decisions/evaluation.md, "Stage attribution";
% backend/src/evaluation/attribution.py (built 20 October 2026).
% ============================================================

\subsection{Stage Attribution}
\label{subsec:harness-attribution}

Each scored claim whose verdict differs from its gold label is
attributed to one stage of the architecture \cite{guo2022survey}, the
first that matches in Table~\ref{tab:attribution}. The rule is applied
twice, because the reference links are the better witness but not
every row has them --- an \texttt{UNVERIFIED} gold label needs none.

\begin{table}[ht]
\centering
\small
\caption{How a wrong verdict is attributed to a stage. The first row
that matches decides.}
\label{tab:attribution}
\begin{tabularx}{\linewidth}{@{}lXX@{}}
\toprule
\textbf{Stage} & \textbf{Reference-based} & \textbf{Trace only} \\
\midrule
Aggregation & The model's raw verdict equals gold and recalibration
changed it & The same \\
Retrieval & No reference link or domain among the candidates & No
candidates at all \\
Ranking & A reference among the candidates, none ranked (cut by the
funnel, the same-domain rule, the pertinence gate or the cap) &
Candidates, but nothing ranked \\
Reasoning & A reference is ranked, and the model's raw verdict differs
from gold & Something ranked, and the model's raw verdict differs from
gold \\
\bottomrule
\end{tabularx}
\end{table}

Both columns are reported, with counts, shares and claim identifiers,
and the claims on which they disagree are listed. Under the reference
rule, a wrong verdict resting on sources other than the annotator's
counts as retrieval: the decisive source was never found. Under the
trace rule the same claim counts as reasoning: the model misread what
it was shown. Both are true statements about the claim, and the
disagreement between them is itself informative. Claims set apart
(search or model unreachable, errors) are never attributed. Everything
is read from fields the records already hold.

Aggregation is tested first. When the model's own verdict was the gold
label, no stage before it caused the error, however little the model was
shown; tested last, a right answer the scorer overrode would be charged
to retrieval whenever the references were missed. An oracle run has no
reference column, since its evidence is the annotator's by
construction.

The rule is a module of the harness since 20~October, a section of
every run's report, tested on records built by hand. For the X-Fact
pilot it was first computed by a script (Section~\ref{sec:results-pilot}),
which now imports the same module.


% ============================================================
% Chapter 6: Experimental Design
%   Section: Evaluation Harness
%     Subsection: Journalist-Effort Metrics
% Included from chapters/06-experimental-design/04-evaluation-harness/evaluation-harness.tex
% Source: docs/decisions/evaluation.md, "Journalist effort" (written
% when the effort question was RQ3; under the questions fixed on
% 14 October it is RQ4). All four built on 20 October 2026:
% labeller/app.py (timing, stats), backend/src/evaluation/effort.py.
% ============================================================

\subsection{Journalist-Effort Metrics}
\label{subsec:harness-effort}

RQ4 asks what share of the journalist's work the system takes on, at
what quality. Four measures answer it.

\begin{itemize}
    \item \textbf{Claim selector's precision}: claims labelled over
    claims labelled plus claims rejected, over every daily batch, in
    total, by rejection reason and by language. The first six facts,
    chosen by hand before the daily batch existed, are excluded.
    Computed by the labelling tool's \texttt{stats} command and the
    evaluation page.
    \item \textbf{Time per fact}: the time from opening a fact's form to
    saving it, counting only the seconds the page is visible, reported
    as median and interquartile range, with the wall time beside it:
    the visible seconds leave out the minutes spent reading sources in
    other tabs, so the truth lies between the two. Recorded by the
    labelling tool since 20~October; facts saved earlier have none.
    \item \textbf{Usable without re-checking}: the share of scored
    claims whose verdict equals the gold label, is definitive, cites at
    least one source with a stance and has neither leak flag. Its
    counterpart is \textbf{confidently wrong}: a definitive verdict
    that differs from gold, which is what a journalist who trusted the
    system would publish. A section of every run's report.
    \item \textbf{Effort estimate}: the usable share times the median
    time per fact, in minutes saved per hundred claims, beside the
    confidently-wrong share as the cost of trusting it --- stated as an
    estimate, with its assumptions, and only when the median time per
    fact is given to the report, which reads no labelling file.
\end{itemize}

For the steps before verification, the same question is asked of
discovery and selection: how many sources and links a scoped round
reads against an every-topic round, how many of the first twenty
places are worth offering, and how far the AI selection agrees with
the editor (Section~\ref{sec:rq-metrics}).

